\section{A finite-scale effective energy with no quadratic width cost}
\label{sec:effective-energy}

Write $H\geq2$ for the number of hidden layers and fix
\[
  \sigma^2=t^2N^{-\gamma},\qquad 0<t\leq1,\qquad\gamma\geq1.
\]
The independent Gaussian weight variances are $N\sigma^2/d$ in the
first layer, $\sigma^2$ in each internal layer, and $\sigma^2/N$ in
the readout.  The likelihood variance is $n\lambda\sigma^2$.  Thus the
posterior is the probability law
\begin{equation}
 \Pi(dw)=Z^{-1}\exp\left[-\frac{\widehat R_S(z_w)}{2\lambda\sigma^2}\right]
 \mu(dw).
 \label{eq:effective-posterior}
\end{equation}
The construction below applies at the displayed finite parameters.
No limit in width or prior strength is used.

The ordinary quadratic energy under this prior has an internal-layer
contribution of order $t^2N^{2-\gamma}$.  This is unbounded in width
when $1\leq\gamma<2$, despite the much smaller operator norm of the
Gaussian matrices.  We retain only their singular values above the
Gaussian operator-norm scale when measuring complexity.  This changes
an auxiliary quantity in the proof; it does not restrict or modify the
learner's prior.

\subsection{Spectral truncation and an exponential moment}

Set $a_N=\sigma\sqrt N=tN^{(1-\gamma)/2}$ and $b_N=8a_N$.
The quantity $a_N$ measures the scale of the Gaussian internal
operator norms.  For each internal
matrix $W^\ell$, let $T^\ell=(W^\ell)_{>b_N}$ retain singular values
strictly greater than $b_N$ and set the others to zero.  Let
$\widetilde w$ be the network with these internal matrices and the
original first layer and readout.  Define
\begin{equation}
 q_N(w):=Q_N^+(\widetilde w)
 =\sum_{\ell=2}^H\|T^\ell\|_F^2+N\|v\|_2^2.
 \label{eq:effective-energy}
\end{equation}
This is a nonnegative Borel function and is at most $Q_N^+(w)$.
The first-layer weights are unrestricted in the capacity estimate
for $Q_N^+$ and are unchanged by truncation, so they need not enter
this auxiliary energy.  They still contribute to the prior and to
the teacher comparison below.

\begin{lemma}[Dimension-free cost of the internal Gaussian bulk]
\label{lem:effective-mgf}
Put
\begin{align}
 C_N&=\frac N2\log\frac87+(H-1)\log\frac32,
 \notag\\
 c_H&=\frac12\log\frac87+(H-1)\log\frac32.
 \label{eq:effective-mgf-constant}
\end{align}
Then
\begin{equation}
 \log\mathbb E_\mu\exp\left(\frac{q_N}{16\sigma^2}\right)
 \leq C_N,
 \qquad \sigma^2 C_N\leq a_N^2c_H.
 \label{eq:effective-mgf}
\end{equation}
Consequently, for every probability law $\nu\ll\mu$ of finite
relative entropy,
\begin{equation}
 \operatorname{KL}(\nu\Vert\mu)
 \geq\frac1{16\sigma^2}\int q_N\,d\nu-C_N.
 \label{eq:effective-entropy}
\end{equation}
\end{lemma}
\begin{proof}
The readout contribution, divided by $\sigma^2$,
has a chi-square distribution with $N$ degrees of freedom.
Its exponential moment at parameter $1/16$ is therefore
$(8/7)^{N/2}$.

For an internal matrix $W=\sigma Z$, where $Z$ has independent standard
Gaussian entries, the Gaussian operator-norm bound
\[
 \mathbb P(\|Z\|_{\mathrm{op}}>2\sqrt N+u)\leq e^{-u^2/2}
\]
implies that the operator-norm ball
$B=\{A:\|A\|_{\mathrm{op}}\leq b_N/2\}$ has probability at least
$1/2$.  For example, $u=2\sqrt N$ gives a failure probability at
most $e^{-2N}<1/2$.  The squared Frobenius distance to this ball is
\[
 \operatorname{dist}_F(W,B)^2
   =\sum_j(s_j(W)-b_N/2)_+^2.
\]
Since $s\leq2(s-b_N/2)$ whenever $s>b_N$,
\[
 \|W_{>b_N}\|_F^2\leq4\operatorname{dist}_F(W,B)^2.
\]
Gaussian isoperimetry and $1-\Phi(u)\leq\tfrac12e^{-u^2/2}$ yield
\begin{equation}
 \mathbb P(\|W_{>b_N}\|_F^2>a)
 \leq\frac12\exp\left(-\frac a{8\sigma^2}\right),
 \qquad a>0.
 \label{eq:effective-internal-tail}
\end{equation}
Tail integration now gives
\begin{align*}
 \mathbb E\exp\left(\frac{\|W_{>b_N}\|_F^2}{16\sigma^2}\right)
 &=1+\frac1{16\sigma^2}\int_0^\infty
      e^{a/(16\sigma^2)}
      \mathbb P(\|W_{>b_N}\|_F^2>a)\,da
 \leq\frac32.
\end{align*}
Independence of the layer blocks proves the first inequality
in~\eqref{eq:effective-mgf}.  Its second inequality follows from $\sigma^2N=a_N^2$
and $\sigma^2\leq a_N^2$.
The entropy variational inequality proves
\eqref{eq:effective-entropy}; truncation of $q_N$ justifies the same
conclusion if its integrability has not yet been established.
\end{proof}

\begin{lemma}[Approximation by the truncated network]
\label{lem:effective-approximation}
For every parameter vector and every input,
\begin{align}
 |z_w(x)-z_{\widetilde w}(x)|
 &\leq b_N\sum_{j=1}^{H-1}q_N(w)^{j/2}
 \notag\\
 &\leq8a_N(H-1)(1+q_N(w))^{(H-1)/2}.
 \label{eq:effective-approximation}
\end{align}
Both outputs have absolute value at most $\sqrt{q_N(w)}$.
\end{lemma}
\begin{proof}
The first hidden layers agree.  Write
$e_\ell=\|h^\ell-\widetilde h^\ell\|_2/\sqrt N$.
Using
\[
 W^\ell h^{\ell-1}-T^\ell\widetilde h^{\ell-1}
 =T^\ell(h^{\ell-1}-\widetilde h^{\ell-1})
      +(W^\ell-T^\ell)h^{\ell-1},
\]
the contraction of tanh, $\|h^{\ell-1}\|_2\leq\sqrt N$, and
$\|T^\ell\|_{\mathrm{op}}\leq\sqrt{q_N}$ give
\[
 e_1=0,\qquad e_\ell\leq\sqrt{q_N}\,e_{\ell-1}+b_N.
\]
Multiplication by $\sqrt N\|v\|_2\leq\sqrt{q_N}$ gives the first
bound.  The second follows from $b_N=8a_N$ and
$q_N^{j/2}\leq(1+q_N)^{(H-1)/2}$.  Bounded hidden activations give
the output bounds.  In particular, no operator-norm contribution
$b_N$ is needed in the multiplying gain of the recursion.
\end{proof}

The posterior variational identity and
Lemma~\ref{lem:effective-mgf} imply, for every comparison law $\nu$
of finite relative entropy,
\begin{align}
 \int\left[\widehat R_S(z_w)+\frac\lambda8q_N(w)\right]\Pi(dw)
 &\leq\int\widehat R_S(z_w)\,\nu(dw)
       +2\lambda\sigma^2\operatorname{KL}(\nu\Vert\mu)
       +2\lambda\sigma^2 C_N.
 \label{eq:effective-variational}
\end{align}
All these quantities are finite for the comparison law constructed
next and for the posterior.  For the latter, its density relative to
$\mu$ is bounded by $1/Z$, and $q_N\leq Q_N$ has all Gaussian moments.

\subsection{A teacher neighborhood with linear width cost}

Let $g$ have width $m$, depth $H+1$, and weights bounded by $B\geq1$.
At any $N\geq m$, balanced cloning computes $g$ exactly and has energy
at most
\begin{equation}
 E=d(d+1)B^2+2Hm^2B^2,
 \qquad |g(x)|\leq U:=mB.
 \label{eq:effective-teacher-energy}
\end{equation}
The neighborhood below narrows only the noise components that act on
the teacher's group-constant activations.  The other Gaussian
coordinates remain present.

\begin{lemma}[Teacher neighborhood at every finite width and scale]
\label{lem:effective-teacher-tube}
For $0<\rho\leq1$, define
\begin{align*}
 b_0&=2\sqrt{2[2\log9+\log(4(H+1))]},
 & A_0&=2+3mB+b_0,\\
 P_0&=1+2mB(d+H+2)A_0^{H-1},
 & \eta&=\min\{1/2,\rho/(2P_0)\},\\
 D_g&=d+2+(H-1)m,
 & c_0&=3/8.
\end{align*}
Let $U_\rho(g)$ be the event that $|z_w(x)-g(x)|<\rho$ at every
$x\in[-1,1]^d$.  Then, for $N\geq m$, $\gamma\geq1$, and $0<t\leq1$,
\begin{equation}
 -\log\mu(U_\rho(g))
 \leq\frac E{2\sigma^2}+ND_g\log\frac1{c_0\eta}+\log2.
 \label{eq:effective-teacher-tube}
\end{equation}
In particular, with $\rho=\lambda/[4(U+1)]$ and $0<\lambda\leq1$,
there is a comparison law $\nu$ with
\begin{align}
 \int\widehat R_S(z_w)\,d\nu
    +2\lambda\sigma^2\operatorname{KL}(\nu\Vert\mu)
 &\leq\widehat R_S(g)+\lambda E+\lambda
 \notag\\
 &\quad+2\lambda a_N^2
       \left[D_g\log\frac1{c_0\eta}+\log2\right].
 \label{eq:effective-comparison}
\end{align}
\end{lemma}
\begin{proof}
Let $w_*$ be the balanced clone and let $P$ project onto its
$m$-dimensional group-constant subspace, with orthonormal basis
$u_j=k_j^{-1/2}\mathbf1_{\text{group }j}$.  In standardized Gaussian
coordinates impose $|\xi|\leq\eta$ on each first-layer and readout
coordinate, and
$|\langle Z_{i,*},u_j\rangle|\leq\eta$ for every internal row and
group.  There are $ND_g$ independent scalar constraints.  A centered
standard Gaussian has probability at least $c_0\eta$ of lying in
$[-\eta,\eta]$ when $0<\eta\leq1$; for example, combine symmetric
unimodality with Chebyshev's probability $3/4$ for $[-2,2]$.
The cylinder therefore has probability at least $(c_0\eta)^{ND_g}$.

Conditioning on this centrally symmetric convex cylinder preserves
centered unit-subgaussian linear forms.  This follows directly by
factorizing into Gaussian coordinates along the orthonormal constrained
subspaces, or from the strongly log-concave moment-generating-function
bound.  Two sphere nets of cardinality at most $9^N$ consequently give
\[
 \mathbb P(\|Z\|_{\mathrm{op}}>2u\sqrt N\mid\text{cylinder})
 \leq2\exp[(2\log9-u^2/2)N].
\]
Intersecting with $\|Z\|_{\mathrm{op}}\leq b_0\sqrt N$ for every
internal layer leaves a symmetric event $A$ with probability at least
$\tfrac12(c_0\eta)^{ND_g}$.

Translate by the standardized clone.  Write
$a_N=\sigma\sqrt N=tN^{(1-\gamma)/2}\leq1$.
The internal perturbation is
$a_NZ/\sqrt N=a_NZP/\sqrt N+a_NZ(I-P)/\sqrt N$.
Its second term annihilates the clone's activations, while the first
has operator norm at most $\eta\sqrt m$.
For the normalized hidden errors this gives, uniformly on the cube,
\[
 e_1\leq\eta(d+1),\qquad
 e_\ell\leq(\sqrt2mB+b_0)e_{\ell-1}+\eta\sqrt m.
\]
Thus $e_H\leq\eta(d+H+2)A_0^{H-1}$.
The clone readout has $\sqrt N\|v_*\|_2\leq\sqrt2mB$, and the
perturbation readout has $\ell^1$ norm at most $\eta$.
The output error is at most $\eta P_0\leq\rho/2$.
Hence the translated event is contained in $U_\rho(g)$.

For standard Gaussian measure $G$ and a centrally symmetric event $A$,
\[
 G(A+a_*)\geq G(A)e^{-\|a_*\|_2^2/2}.
\]
Indeed, the Gaussian density ratio averaged over $A$ is
$e^{-\|a_*\|_2^2/2}\mathbb E[e^{-\langle Z,a_*\rangle}\mid A]$,
and its second factor is at least one by symmetry and Jensen.
Here $\|a_*\|_2^2=Q_N(w_*)/\sigma^2\leq E/\sigma^2$.
This proves~\eqref{eq:effective-teacher-tube}.

For the last assertion, use $\nu=\mu(\cdot\mid U_\rho(g))$.
Its relative entropy is $-\log\mu(U_\rho(g))$, and
\[
 \int\widehat R_S(z_w)\,d\nu
 \leq\widehat R_S(g)+2\rho(U+1)+\rho^2
 \leq\widehat R_S(g)+\lambda.
\]
Finally $\sigma^2N=a_N^2$ and
$\sigma^2\leq a_N^2$ give~\eqref{eq:effective-comparison}.
\end{proof}
